% ============================================================
% header.tex  --  aturan tipografi sesuai Panduan TA/Skripsi DTNTF UGM
% ============================================================
\usepackage{titlesec}
\usepackage{titletoc}
\usepackage{indentfirst}
\usepackage{fancyhdr}
\usepackage{caption}
\usepackage{etoolbox}
\usepackage{setspace}

% Alinea rata kiri-kanan, awal alinea menjorok 1 cm
\setlength{\parindent}{1cm}
\setlength{\parskip}{0pt}

% ---------- Penomoran ----------
% Bab: BAB I (Romawi); subbab: I.1, I.2.1
% Gambar/Tabel/Persamaan: angka Arab per bab (3.1). Counter di-reset per bab oleh class report.
\renewcommand{\thechapter}{BAB~\Roman{chapter}}
\renewcommand{\thesection}{\Roman{chapter}.\arabic{section}}
\renewcommand{\thesubsection}{\thesection.\arabic{subsection}}
\renewcommand{\thesubsubsection}{\thesubsection.\arabic{subsubsection}}
\renewcommand{\thefigure}{\arabic{chapter}.\arabic{figure}}
\renewcommand{\thetable}{\arabic{chapter}.\arabic{table}}
\renewcommand{\theequation}{\arabic{chapter}.\arabic{equation}}

% ---------- Judul bab / subbab ----------
% Judul bab: HURUF KAPITAL, tebal, rata tengah, tanpa titik
\titleformat{\chapter}[display]
  {\normalfont\bfseries\centering}{\thechapter}{12pt}{\MakeUppercase}
\titleformat{name=\chapter,numberless}[block]
  {\normalfont\bfseries\centering}{}{0pt}{\MakeUppercase}
\titlespacing*{\chapter}{0pt}{0pt}{24pt}

% Subbab dst.: rata kiri, tebal, Title Case (diketik manual), tanpa titik di akhir
\titleformat{\section}{\normalfont\bfseries}{\thesection.}{0.6em}{}
\titleformat{\subsection}{\normalfont\bfseries}{\thesubsection.}{0.6em}{}
\titleformat{\subsubsection}{\normalfont\bfseries}{\thesubsubsection.}{0.6em}{}
\titlespacing*{\section}{0pt}{14pt}{6pt}
\titlespacing*{\subsection}{0pt}{12pt}{6pt}
\titlespacing*{\subsubsection}{0pt}{10pt}{6pt}

% ---------- Daftar isi ----------
\titlecontents{chapter}[0pt]{\addvspace{6pt}}
  {\thecontentslabel\ \ }{}{\titlerule*[.5pc]{.}\contentspage}
\titlecontents{section}[4em]{}
  {\contentslabel{3em}}{}{\titlerule*[.5pc]{.}\contentspage}
\titlecontents{subsection}[8em]{}
  {\contentslabel{4em}}{}{\titlerule*[.5pc]{.}\contentspage}

% ---------- Nomor halaman ----------
% Kanan atas; tengah bawah pada halaman yang memuat judul bab
\pagestyle{fancy}
\fancyhf{}
\fancyhead[R]{\thepage}
\renewcommand{\headrulewidth}{0pt}
\fancypagestyle{plain}{%
  \fancyhf{}\fancyfoot[C]{\thepage}\renewcommand{\headrulewidth}{0pt}}

% ---------- Judul gambar/tabel ----------
% "Gambar 3.1." / "Tabel 3.1." (label tebal), spasi 1 untuk judul >1 baris
\captionsetup{font={stretch=1}, labelfont=bf, labelsep=period,
              singlelinecheck=off, justification=justified, format=hang}
\DeclareCaptionListFormat{gambarlist}{\figurename~#1#2}
\DeclareCaptionListFormat{tabellist}{\tablename~#1#2}
\captionsetup[figure]{listformat=gambarlist}
\captionsetup[table]{listformat=tabellist}
\makeatletter
\renewcommand*{\l@figure}{\@dottedtocline{1}{0em}{6em}}
\renewcommand*{\l@table}{\@dottedtocline{1}{0em}{6em}}
\makeatother

% ---------- Daftar pustaka: spasi 1 ----------
\AtBeginEnvironment{CSLReferences}{\singlespacing}
